\question{}

After the success of the kernel you developed in ``Question 3'',
you now want to target applications with much larger $n$.
For this reason, you now want to solve it with several compute nodes.
You have access to $m$ compute nodes and enough material to connect $m$ pairs of compute nodes together.
Using these connections, the communication of $k$ bytes takes a time
equal to $\alpha + \beta k$.

Answer the following questions:
\begin{parts}
  \part How would you interconnect the $m$ compute nodes ?
    \begin{solutionbox}{7cm}
      We could use $n/m$ adjacent values for each node and then connect each node with the two nodes
      holding adjacent values as a ring.
      A tree would also be possible, it would only require connecting $m-1$ pairs but some neighbors
      won't be connected directly in a tree and will need $\log_2(m)$ communications.
    \end{solutionbox}
  \part How much time would the communication take at each time step of the integration
    in terms of $\alpha, \beta, m, n$ ?
    \begin{solutionbox}{7cm}
      We can first communicate the values of \lstinline|prev| and then the values of \lstinline|next|.
      That's two communications, each of 4 bytes:
      \[ 2\alpha + 8\beta. \]
    \end{solutionbox}
  \part You would like to compute $\sum_{i=1}^n |u(\theta_i, t)|$
    to verify whether the solution is diverging.
    Which MPI collectives would you use ?
    How much time would this take in terms of
    $\alpha, \beta, m, n$ and also $\gamma$, the
    time taken by the compute node to execute one arithmetic operation.
    \begin{solutionbox}{23cm}
      We can use \lstinline|MPI_reduce| with \lstinline|MPI_sum|.
      We cannot use the classical complexity of \lstinline|MPI_reduce| with $\log_2(m)$ since we are using a
      ring. As the diameter is $m/2$, we need at least $m/2$ communications.
      Each communication only needs to communicate the value of the sum so it is only 4 bytes.
      The complexity of doing the sum in each node is $2n/m \gamma$ (we multiply by 2 since we both need to do \lstinline|+| and the absolute value).
      One could also consider that each node had access to at least $n/m$ threads in which case the complexity drops down to $\log_2(n/m) \gamma$,
      this was also accepted.
      We also need to merge the partial sums. This case the complexity depends on the approach, both were accepted (since this part of the complexity would hardly be the bottleneck):
      Let $x_1, \ldots, x_m$ be the partial sums.
      \begin{itemize}
        \item We could
        \begin{enumerate}
          \item first sum $x_1, x_2$ into $x_{2}$ on node 2 and sum $x_{m-1},x_m$ into $x_{m_1}$ on node $m-1$.
          \item Next we merge $x_2,x_3$ into $x_3$ on node 3 and $x_{m-1},x_{m-2}$ on node $m-2$,
          \item etc...
        \end{enumerate}
        The complexity will then be $m/2 \gamma$.
        \item We could
        \begin{enumerate}
          \item sum $x_{2i}, x_{2i-1}$ into $x_{2i}$ on node $2i$ for $i = 1, \ldots, m/2$.
          \item Then we communicate $x_{4i-2}$ to $x_{4i-1}$ (no arithmetic operation needed) on node $4i-1$ for $i = 1, \ldots, m/4$.
          \item Then we sum $x_{4i-1}, x_{4i}$ into $x_{4i}$ on node $4i$ for $i = 1, \ldots, m/4$.
          \item Then we communicate $x_{8i-4}$ to $x_{8i-3}$ (no arithmetic operation needed) on node $8i-3$ for $i = 1, \ldots, m/8$.
          \item Then we communicate $x_{8i-3}$ to $x_{8i-2}$ (no arithmetic operation needed) on node $8i-2$ for $i = 1, \ldots, m/8$.
          \item Then we communicate $x_{8i-2}$ to $x_{8i-1}$ (no arithmetic operation needed) on node $8i-1$ for $i = 1, \ldots, m/8$.
          \item etc...
        \end{enumerate}
        This would require more communication but since they are in parallel it gives the same amount of communication time (but this might consume more energy though).
        The complexity in terms of $\gamma$ is decreased to $\log_2(m)\gamma$.
      \end{itemize}

      In total, we have
      \[ m/2 (\alpha + 4 \beta + \gamma) + 2n/m \gamma. \]
      As it is a complexity, constants didn't matter so, the following was also accepted for instance:
      \[ m (\alpha + \beta + \gamma) + n/m \gamma. \]
    \end{solutionbox}
\end{parts}
